\begin{center}
{\Large\bfseries 基于体检指标与异质模型融合的\\[0.25em]血糖回归预测及预处理对照研究}
\end{center}

\begin{abstract}
针对利用临床及体检指标预测个体血糖浓度的竞赛任务，本文在保留全部待评价记录的条件下，研究缺失处理、特征构造、模型选择与异质回归模型融合。将原训练集与带标签A集合并，固定划分为4649条训练记录和1993条开发记录；缺失填补、标准化与可选主成分分析均限定在相应训练部分拟合。首先比较六类基准模型，随后依次进行特征分组消融、有限预算联合参数搜索、树与线性模型融合、小型神经网络及其补充融合，利用训练折外预测确定配置与权重。最终方案由优化XGBoost、LightGBM、Ridge和去除乙肝五项的小网络各占25\%组成，其开发集和B集均方误差分别为2.177243和1.234296，相对原XGBoost分别降低约2.31\%和1.47\%。进一步固定十个模型版本的特征、参数和网络轮数，开展标准化与主成分分析的四处理对照。结果表明，本轮两个网络从标准化中获得较低的训练折外误差，额外增加主成分分析未显示稳定收益；仅调整推荐预处理的融合版本在开发集略有改善，在B集略有变差。最终融合对高血糖值仍有低估，B集相对基线的配对误差差异区间跨零。由于开发集和B集均有历史研究使用，结果属于已知数据上的复核，尚不能证明未来样本上的稳定优势。
\end{abstract}
\noindent\textbf{关键词：}血糖预测；体检数据；缺失值处理；模型融合；预处理对照
